%% sections/conclusion.tex — Step B12
\section{Conclusions}
\label{sec:conclusion}

We extended the DyAM multimodal attention framework for ICI response
prediction in NSCLC with a competitive (OvO) attention mechanism and an
NLP-based encoding of structured clinical variables.
Under a repeated-seed evaluation protocol, OvO attention matched the
original cooperative attention at every primary benchmark configuration
(within 0.011~AUC, $p \geq 0.17$; AUC~$=0.773$ vs.\ $0.760$ at the primary
benchmark), offering a simpler, $N$-parameter alternative to the
$N \times N$ cooperative weight matrix without a measurable accuracy cost;
an initial single-partition screen had suggested a larger OvO advantage in
low-modality configurations (best single-partition AUC~$=0.800$, later
$0.8133$ when combined with NLP-clinical encoding), and averaging across
five independent CV partitions confirmed OvO's accuracy while showing this
particular single-partition margin does not reproduce exactly.
NLP encoding of clinical features using a general-purpose sentence
transformer consistently improved AUC over raw numeric encoding
(up to $\Delta\text{AUC} = +0.030$) and provided the strongest independent
prognostic value in survival analysis, with the highest Cox hazard ratios and
time-dependent AUC among all clinical encoding strategies, despite AUC
differences not reaching statistical significance at $n=247$.
This improvement was confirmed under repeated-seed evaluation (NLP
encoding exceeded numeric labs in 19 of 21 modality combinations) and
generalised across fusion mechanisms: combining NLP-clinical encoding with
OvO attention reproduced the same fusion-dependent gain, indicating that
the clinical-encoding strategy, rather than the choice of attention
mechanism, drives the improvement.
These findings suggest that sentence-embedding-based representations of
tabular clinical data are a practical and generalisable strategy for
multimodal biomarker discovery, robust to the underlying attention
architecture.
Validation in larger, multi-institutional cohorts is needed to confirm these
findings and to establish their clinical utility for treatment selection.
